一份实验报告写着：主指标 \(p=0.31\)，判定无效；十四个指标里有两个显著，讲了个故事；结论栏那句``新方案有效的概率为 \(97\%\)''来自另一次实验的 \(p=0.03\)。三句话，三处错，而且都是常见错。这一单元处理的就是这三处。

假设检验是一条在看数据之前写定的决策规则。规则会犯两类错误：原假设为真却拒绝，与备择为真却没拒绝。前者用 \(\alpha\) 一次性锁死，后者随效应大小变化，它的补是功效。样本量、效应量、显著性水平、功效被一个等式拴住，知道三个就能算第四个——实验设计的算术到此为止。

规则一旦离开``事先固定''这个前提，所有承诺同时失效。看到显著就停，同时看十几个指标只报最好的那个，事后挑单侧方向，都属于让数据参与了规则的制定。举一个数字：\(20\) 个相互独立且原假设全真的检验，每个用 \(0.05\) 的水平，至少误报一次的概率是 \(1-(1-0.05)^{20}\approx 0.64\)。计算没有变，计量单位从``每次检验''变成了``整族检验''。

本单元按决策逻辑推进：

\begin{enumerate}
\tightlist
\item
  两类错误与功效共同定义检验——把检验当成决策规则，讲清临界值如何在两类错误之间搬运，以及四个量的联动怎样支撑样本量估算。
\item
  \(p\) 值不是原假设为真的概率——从尾概率的定义出发纠正方向混淆，逐条拆解常见误读，说明可选停止与隐藏选择为什么会让 \(p\) 值失去校准。
\item
  似然比比较两个模型对数据的解释力——用固定预算下的装箱直觉给出 Neyman--Pearson 引理，再走到广义似然比、Wilks 近似及其失效场合。
\item
  多重比较改变错误的计量单位——区分族错误率与错误发现率，讲清 Bonferroni 与 Benjamini--Hochberg 各自适合什么决策形状。
\end{enumerate}

检验交付的是特定模型、特定设计、特定错误控制之下的一个决定。它不等于效应大小，不等于实际重要性，也不等于因果结论。报告里应当同时出现效应估计与区间，把不确定性的宽度交给读者。

\subsection{怎么读这一章}

四篇按依赖顺序排列，先读第一篇建立框架，第二篇是本单元误解最集中的地方，值得慢读。若手头正有一份实验报告或一次特征筛选要处理，也可以直接翻到最接近当前困惑的一篇，再沿篇内链接回补。读每一篇时反复追问同一件事：原假设是什么，拒绝规则是什么，宣称的错误率针对的是哪一个重复过程。

\subsection{本章篇目}

以下篇目按概念依赖排列：

\begin{itemize}
\tightlist
\item \hyperref[sec:11-Mathematical-Statistics/05-Hypothesis-Testing/01]{``两类错误与功效共同定义检验''}；
\item \hyperref[sec:11-Mathematical-Statistics/05-Hypothesis-Testing/02]{``\texorpdfstring{\(p\) 值不是原假设为真的概率}{p 值不是原假设为真的概率}''}；
\item \hyperref[sec:11-Mathematical-Statistics/05-Hypothesis-Testing/03]{``似然比比较两个模型对数据的解释力''}；
\item \hyperref[sec:11-Mathematical-Statistics/05-Hypothesis-Testing/04]{``多重比较改变错误的计量单位''}。
\end{itemize}
